\chapter{Durchführung und Methodik der Datenerhebung}
\label{chap:methodik}

Während das vorangegangene Kapitel das Forschungsdesign und die Hypothesen definiert hat, beschreibt dieses Kapitel den konkreten Ablauf der Studie sowie die Operationalisierung der Messgrößen. Es wird detailliert dargelegt, wie die Datenerhebung organisiert wurde, um die interne Validität der Ergebnisse sicherzustellen.

\section{Ablauf der Untersuchung}

Die Studie wurde als kontrolliertes Experiment durchgeführt. Jede Sitzung folgte einem strikten Protokoll, um Störfaktoren zu minimieren. Der Ablauf gliederte sich in fünf Phasen:

\begin{enumerate}
    \item  \textbf{Onboarding (10 Min.):}
    Die Teilnehmenden erhielten eine Einführung in das Szenario (\enquote{Sie sind neuer Entwickler im Team...}) und die Entwicklungsumgebung. Es wurde sichergestellt, dass GitHub Copilot aktiviert und funktionsfähig war.
    
    \item  \textbf{Pre-Survey (5 Min.):}
    Erfassung der demografischen Daten und Selbsteinschätzung zur Kompetenzeinstufung (siehe Kapitel~\ref{chap:forschungsdesign}).
    
    \item  \textbf{Vibe-Coding-Challenge (60 Min.):}
    Die Kernphase der Untersuchung. Die Teilnehmenden bearbeiteten die Aufgaben selbstständig. Eingriffe durch den Versuchsleiter erfolgten nur bei technischen Problemen (z.\,B. Absturz der IDE), nicht jedoch bei inhaltlichen Fragen.
	\csname textbf\endcsname{Begründung des Zeitlimits:} Die Begrenzung auf 60 Minuten dient der organisatorischen Machbarkeit und fungiert als Stressor. Zeitdruck erhöht die Wahrscheinlichkeit, dass Teilnehmende auf heuristische Strategien (\emph{System 1}) zurückgreifen. Dies kann den \emph{Automation Bias} verstärken und simuliert realistische Deadline-Situationen.
    
    \item  \textbf{Datensicherung (5 Min.):}
    Sicherung der Arbeitsartefakte (Branch-Stand, Test-/Lint-Ausgaben und Chat-Export) als Grundlage der späteren Artefaktanalyse.
    
    \item  \textbf{Debriefing (5 Min.):}
    Kurzes Abschlussgespräch, um offene Fragen zu klären und qualitatives Feedback einzuholen, das im Survey nicht abgebildet werden konnte.
\end{enumerate}

\section{Operationalisierung der Metriken}

Um die Hypothesen zu prüfen, wurden die abstrakten Beobachtungsdimensionen (Aufgabenfortschritt/Validierung sowie Interaktion) in nachvollziehbare, artefaktgebundene Indikatoren übersetzt.

\subsection{Aufgabenfortschritt und Validierung (Pipeline-Status)}
Die Ergebnislogik der Arbeit basiert auf einer artefaktgestützten Pipeline-Klassifikation pro Task und Branch. Zur Trennung von Intention, Implementationsartefakt und testgebundener Bestätigung werden die Statuskategorien  \textit{Attempted},  \textit{Implemented\textsubscript{static}} und  \textit{Verified} verwendet (Definitionen und Trichter-Evidenzstufe  \textit{Implemented$^\ast$} in Abschnitt~\ref{sec:results_pipeline}).
Wichtig ist dabei:  \textit{Verified} ist strikt testgebunden (\enquote{besteht die aktuellen Tests}) und somit ein methodengebundenes Diagnose-Signal, kein allgemeines Qualitätsurteil.

\subsection{Code- und Test-Artefakte (Surrogatindikatoren)}
Zur Einordnung der Umsetzungsstände werden zusätzlich struktur- und testnahe Artefakte herangezogen. Dazu zählen u.a. heuristische Strukturindikatoren und aggregierte Testausgänge aus den Auswertungsartefakten (z.\,B. \texttt{evaluation/deep\_code\_analysis.json}, \texttt{evaluation/code\_quality\_metrics.json}). Diese Größen werden durchgängig als Surrogate berichtet und nicht als vollständige Qualitätsmessung verstanden.

\subsection{Chat-Artefakte (Interaktion im Vibe-Coding)}
Die Interaktion wird über Chat-Exports und daraus abgeleitete deskriptive Indikatoren beschrieben (z.\,B. Wortanzahl, Prompt-Antwort-Zyklen, Fehlermarker-Iterationen). Die Auswertung bleibt auf Musterbeschreibung und Triangulation mit Pipeline- und Code/Test-Artefakten beschränkt; Chat-Intensität wird nicht als Effizienz- oder Qualitätsmaß interpretiert.

\section{Güte der Untersuchung (Threats to Validity)}

Die Validität der Studie wird anhand der Kategorien von \textcite{cook1979quasi} diskutiert.

\subsection{Konstruktvalidität (Construct Validity)}
Messen wir wirklich das, was wir messen wollen?
\begin{itemize}
    \item  \textbf{Problem:} \enquote{Kompetenz} ist schwer zu quantifizieren.
    \item  \textbf{Mitigation:} Kompetenz wird als Analyseperspektive über eine auditierbare, diskrete Pre-Survey-Heuristik operationalisiert (Tabelle~\ref{tab:competence-rubric}). Es wird explizit keine psychometrische Validierung beansprucht; die Einordnung wird über Triangulation mit Artefaktspuren (Pipeline/Code/Test/Chat) abgesichert.
\end{itemize}

\subsection{Interne Validität (Internal Validity)}
Können die beobachteten Effekte eindeutig der unabhängigen Variable (Kompetenz) zugeschrieben werden?
\begin{itemize}
    \item  \textbf{Selection Bias:} Die Zuteilung zu den Gruppen basierte auf Selbsteinschätzung. Um dies zu mitigieren, wurde der Algorithmus zur Klassifizierung strikt angewendet.
    \item  \textbf{Experimenter Bias (Rosenthal-Effekt):} Die Erwartungshaltung des Versuchsleiters könnte die Teilnehmenden beeinflussen.  \textbf{Gegenmaßnahme:} Standardisierte Instruktionen und minimale Interaktion während der Coding-Phase.
    \item  \textbf{Compensatory Rivalry (John Henry Effect):} Wenn eine Gruppe sich benachteiligt fühlt, könnte sie sich mehr anstrengen. Da alle Gruppen Copilot nutzten, ist dieses Risiko gering. Dennoch könnten Teilnehmende versuchen, besonders \enquote{sorgfältig} oder \enquote{kontrolliert} zu arbeiten (z.\,B. durch mehr manuelles Prüfen), was den beobachteten Fortschritt unter Zeitlimit beeinflussen kann \parencite{saretsky1972johnhenry}.
\end{itemize}

\subsection{Externe Validität (External Validity)}
Sind die Ergebnisse verallgemeinerbar?
\begin{itemize}
    \item  \textbf{Hawthorne-Effekt:} Das Wissen, beobachtet zu werden, verändert das Verhalten. Dies ist in Laborstudien typischerweise nicht vollständig vermeidbar, wurde aber durch die realistische Aufgabenstellung (Flow-Erleben) abgemildert.
    \item  \textbf{Ecological Validity:} Die Nutzung einer Brownfield-Umgebung erhöht die Übertragbarkeit auf Unternehmenskontexte im Vergleich zu synthetischen LeetCode-Aufgaben deutlich.
\end{itemize}

\subsection{Reliabilität (Reliability)}
Sind die Ergebnisse reproduzierbar?
\begin{itemize}
    \item  \textbf{Reproduzierbarkeit der Auswertung:}\\
    Die Pipeline-Status und Artefakt-Surrogate werden aus versionierten Branches und gespeicherten Auswertungsartefakten abgeleitet.
    Die Klassifikation folgt festen, dokumentierten Regeln (Attempted, Implemented\textsubscript{static} und Verified; Abschnitt~\ref{sec:results_pipeline}) und ist damit prinzipiell wiederholbar.
\end{itemize}

\section{Datenanalyse-Strategie}

Die Auswertung folgt einem artefaktbasierten, überwiegend deskriptiven Mixed-Methods-Ansatz:

\begin{enumerate}
    \item  \textbf{Deskriptive Quantifizierung:} Pipeline-Status (Attempted, Implemented\textsubscript{static} und Verified; sowie Implemented$^\ast$ für die Trichterdarstellung) werden pro Task und Branch berichtet. Ergänzend werden aggregierte Surrogatindikatoren aus Code- und Test-Artefakten ausgewiesen.
    \item  \textbf{Musteranalyse \& Triangulation:} Chat- und Strukturmuster werden als qualitative Befundlinien beschrieben und mit Pipeline- und Artefaktspuren trianguliert, um Integrations- und Validierungsengpässe im Brownfield-Kontext einzuordnen.
\end{enumerate}

Dieser methodische Aufbau stellt sicher, dass nicht nur das \emph{Ergebnis} (Code), sondern auch der \emph{Prozess} (Interaktion) beleuchtet wird.
